% SPDX-FileCopyrightText: Copyright (c) 2026 The Newton Developers
% SPDX-License-Identifier: CC-BY-4.0

\documentclass[11pt]{article}

\input{glyphtounicode}
\pdfgentounicode=1
\usepackage[T1]{fontenc}
\usepackage{lmodern}
\usepackage[margin=1in]{geometry}
\usepackage{amsmath}
\usepackage{amssymb}
\usepackage{array}
\usepackage{graphicx}
\usepackage[colorlinks=true,linkcolor=blue,citecolor=blue,urlcolor=blue]{hyperref}

\newcommand{\R}{\mathbb{R}}
\newcommand{\dd}{\Delta}
\newcommand{\eps}{\varepsilon}
\newcommand{\pinv}{M^{-1}}
\newcommand{\T}{\mathsf{T}}
\newcommand{\proj}{\Pi}
\newcommand{\clamp}{\operatorname{clamp}}

\title{A Split Stretch/Slip Formulation for Routed Cable Joints}
\author{Newton Developers}
\date{May 2026}
\hypersetup{
  pdftitle={A Split Stretch/Slip Formulation for Routed Cable Joints},
  pdfauthor={Newton Developers}
}

\begin{document}
\maketitle

\begin{abstract}
This note describes the reduced routed-cable model used by Newton's tendon
solver for massless cables over attachments, zero-radius pinholes, and
finite-radius rolling pulleys.  The formulation separates the cable solve into
three local mechanisms: unilateral free-span stretch, capstan slip represented
as rest-length redistribution, and a capstan-limited rolling spin row.  The
same rows are used by the XPBD solver and by the VBD rigid-body path.  This
document states the equations of motion, the exact constraint rows, and the
rest-length projection derived directly from the Euler--Eytelwein capstan
equation.
\end{abstract}

\section{Model}

Let a tendon be an ordered sequence of links $i=0,\ldots,n-1$.  Consecutive
links define free spans $s=0,\ldots,n-2$.  Each link is attached to a rigid
body $b_i$ and has one of three types:
\[
  \texttt{ATTACHMENT}, \qquad \texttt{PINHOLE}, \qquad \texttt{ROLLING}.
\]
Each free span stores a mutable rest length $r_s$, compliance $c_s$, damping
$d_s$, and stretch multiplier $\lambda_s$.  The rest-length distribution is the
state that records how much cable material is assigned to each free span.

For body $b$, let
\[
  q_b = (p_b,Q_b), \qquad v_b = (u_b,\omega_b),
\]
where $p_b$ is translation, $Q_b\in SO(3)$ is orientation, $u_b$ is linear
velocity, and $\omega_b$ is angular velocity.  A body-local point $o$ has world
position
\[
  x_b(o) = p_b + Q_b o.
\]
Let $c_b^0$ be the local center of mass and $x_b^c=p_b+Q_b c_b^0$ its world
position.  The velocity of $x=x_b(o)$ is
\[
  \dot{x} = u_b + \omega_b \times (x-x_b^c).
\]

\section{Equations of Motion and Discrete Projection}

The unconstrained rigid-body dynamics can be written in generalized velocity
form as
\begin{equation}
  M(q)\dot{v} = f_\mathrm{ext}(q,v,t) + J(q,r)^\T \tau,
  \label{eq:eom}
\end{equation}
where $M$ is block-diagonal in body space, $J$ contains the active cable rows,
$r$ is the vector of free-span rest lengths, and $\tau$ are row impulses or
forces depending on time discretization.

Newton applies the cable model as local position/velocity projections.  After
external integration predicts $(q^\star,v^\star)$, each active row computes an
increment $\dd\lambda$ and applies the body-space impulse
\[
  \dd p_b = J_{b,v}^\T\dd\lambda,\qquad
  \dd \ell_b = J_{b,\omega}^\T\dd\lambda.
\]
The velocity increment is
\[
  \dd u_b = m_b^{-1}\dd p_b,\qquad
  \dd\omega_b = I_b^{-1}\dd\ell_b,
\]
using the same effective inverse mass and inverse inertia arrays as the solver
path.  The pose is advanced by the velocity increment over the substep $h$.
For rigid-limit rows $c_s\to 0$, the stretch constraints satisfy the unilateral
complementarity conditions
\begin{equation}
  C_s(q,r_s)\le 0,\qquad \lambda_s\ge 0,\qquad
  \lambda_s C_s(q,r_s)=0.
\end{equation}
Finite compliance replaces this hard complementarity with the XPBD row below.

\section{Route Update}

For an attachment or pinhole, the route point is simply the body-local point
transformed to world space.  For a rolling link with center $g$, radius $R$,
axis $a$, and orientation sign $\sigma\in\{-1,+1\}$, the tangent point from an
external point $p$ is computed in the plane orthogonal to $a$.

Let
\[
  d = g-p,\qquad
  \bar d = d - (d\cdot a)a,\qquad
  u = \frac{\bar d}{\|\bar d\|},\qquad
  v = a\times u.
\]
For $\|\bar d\|>R$, define
\[
  \phi = \arcsin\!\left(\min\left(\frac{R}{\|\bar d\|},1\right)\right),
  \qquad
  \psi =
  \begin{cases}
    -\pi/2-\phi, & \sigma>0,\\
    \phantom{-}\pi/2+\phi, & \sigma<0.
  \end{cases}
\]
The tangent point is
\begin{equation}
  x_\mathrm{tan}=g+R(\cos\psi\,u+\sin\psi\,v).
\end{equation}
For a span between two rolling links, the two tangent points are iterated to
consistency.  A pure spin about the circular pulley axis does not change these
tangent points; this is why the stretch row removes the spin-axis angular
component.

When a rolling contact point moves from $x_\mathrm{old}$ to $x_\mathrm{new}$ on
the pulley surface, the signed surface transfer used by the route update is
\begin{equation}
  \dd \ell
  = R\sigma\,
    \operatorname{atan2}\!\left(
      ((u_\mathrm{new}\times u_\mathrm{old})\cdot a),
      u_\mathrm{old}\cdot u_\mathrm{new}
    \right),
  \label{eq:surface-transfer}
\end{equation}
where $u_\mathrm{old}$ and $u_\mathrm{new}$ are the normalized radial vectors
from the pulley center, projected into the pulley plane.  The signed transfer
is applied to the adjacent span rest length with the capstan factor
\begin{equation}
  \alpha = \exp(\clamp(\mu\theta,0,20)),\qquad
  \beta = \frac{\alpha-1}{\alpha+1},
  \label{eq:alpha-beta}
\end{equation}
so that $\mu=0$ gives no surface drag and high friction approaches the no-slip
surface-transfer limit.

\begin{figure}[t]
  \centering
  \includegraphics[width=0.88\linewidth]{docs/images/cable_joints_formulation/capstan_geometry.pdf}
  \caption{A single capstan row.  The circular route update supplies the
  tangent points and wrap angle $\theta$.  The stretch rows carry the adjacent
  span tensions $T_l,T_r$, while slip changes only the rest-length distribution
  subject to $r_l+r_r=\mathrm{const}$ and the capstan band.}
  \label{fig:capstan-row}
\end{figure}

\section{Free-Span Stretch Row}

Span $s$ runs from left endpoint $x_s^L$ on body $b_L$ to right endpoint
$x_s^R$ on body $b_R$.  Its geometric length, unit direction, and unilateral
stretch constraint are
\begin{equation}
  L_s = \|x_s^R-x_s^L\|,\qquad
  n_s = \frac{x_s^R-x_s^L}{\max(L_s,\eps)},\qquad
  C_s(q,r_s)=L_s-r_s \le 0.
  \label{eq:stretch-constraint}
\end{equation}
If $C_s\le 0$, the span is slack: the row is inactive, $\lambda_s$ is reset to
zero, and the span contributes no tension.

For an active taut span, the translational Jacobian blocks are
\begin{equation}
  J_{L,v}=-n_s,\qquad J_{R,v}=n_s.
\end{equation}
Let
\[
  \rho_L=x_s^L-x_{b_L}^c,\qquad \rho_R=x_s^R-x_{b_R}^c.
\]
The unprojected angular blocks are
\begin{equation}
  \widehat J_{L,\omega}= -\rho_L\times n_s,\qquad
  \widehat J_{R,\omega}= \phantom{-}\rho_R\times n_s.
  \label{eq:unprojected-angular}
\end{equation}
For an endpoint that lies on a rolling link with world axis $a$, the stretch
row removes the spin-axis component:
\begin{equation}
  J_{\omega} =
  \proj_a^\perp \widehat J_{\omega},\qquad
  \proj_a^\perp = I-aa^\T.
  \label{eq:spin-projection}
\end{equation}
For attachments and pinholes, $J_{\omega}=\widehat J_{\omega}$.  Equation
\eqref{eq:spin-projection} is a modeling choice, not a numerical tweak: pulley
spin does not change the route tangent geometry, so the stretch row should not
accidentally impose a no-slip pulley constraint.

The row velocity and effective inverse mass are
\begin{align}
  \dot C_s &=
  J_{L,v}\cdot u_L + J_{L,\omega}\cdot \omega_L
  + J_{R,v}\cdot u_R + J_{R,\omega}\cdot \omega_R,
  \\
  w_s &= J_s M^{-1}J_s^\T.
\end{align}
Written in body blocks,
\begin{equation}
  w_s =
  m_L^{-1}\|J_{L,v}\|^2
  + J_{L,\omega}^\T I_L^{-1}J_{L,\omega}
  + m_R^{-1}\|J_{R,v}\|^2
  + J_{R,\omega}^\T I_R^{-1}J_{R,\omega}.
\end{equation}

Newton's implemented XPBD increment is
\begin{equation}
  \gamma_s=c_s d_s,\qquad
  \dd\lambda_s =
  -\frac{
    C_s + c_s\lambda_s + \gamma_s\dot C_s
  }{
    (h+\gamma_s)w_s + c_s/h
  }.
  \label{eq:xpbd-stretch}
\end{equation}
The solver stores both $\lambda_s\leftarrow\lambda_s+\dd\lambda_s$ and the
latest $\dd\lambda_s$.  The latter is used by the rolling spin row.

The stretch row applies the impulse-like body deltas
\begin{equation}
  \dd P_{b_L}=J_{L,v}^\T\dd\lambda_s,\quad
  \dd L_{b_L}=J_{L,\omega}^\T\dd\lambda_s,\quad
  \dd P_{b_R}=J_{R,v}^\T\dd\lambda_s,\quad
  \dd L_{b_R}=J_{R,\omega}^\T\dd\lambda_s.
  \label{eq:stretch-deltas}
\end{equation}
Relaxation multiplies these deltas before they are accumulated.

\section{Capstan Slip as Rest-Length Projection}

For a guide $g$ between a left span $l$ and right span $r$, define taut
extensions and tension estimates
\begin{equation}
  e_l=\max(L_l-r_l,0),\qquad
  e_r=\max(L_r-r_r,0),\qquad
  T_l=\frac{e_l}{\max(c_l,\eps)},\qquad
  T_r=\frac{e_r}{\max(c_r,\eps)}.
  \label{eq:tension-estimate}
\end{equation}

For a pinhole, the bend angle is
\begin{equation}
  \theta =
  \operatorname{atan2}\!\left(
    \|d_\mathrm{in}\times d_\mathrm{out}\|,
    d_\mathrm{in}\cdot d_\mathrm{out}
  \right),
  \label{eq:pinhole-angle}
\end{equation}
where $d_\mathrm{in}$ is the incoming unit cable direction into the pinhole and
$d_\mathrm{out}$ is the outgoing unit direction.  For a rolling link,
\begin{equation}
  \theta =
  \left|
  \operatorname{atan2}\!\left(
    (u_l\times u_r)\cdot a,\; u_l\cdot u_r
  \right)
  \right|,
  \label{eq:rolling-angle}
\end{equation}
where $u_l$ and $u_r$ are the normalized radial vectors from the pulley center
to the two tangent points, projected into the pulley plane.

The Euler--Eytelwein capstan law gives the admissible static tension band
\begin{equation}
  T_l \le \alpha T_r,\qquad
  T_r \le \alpha T_l,\qquad
  \alpha=\exp(\clamp(\mu\theta,0,20)).
  \label{eq:capstan-band}
\end{equation}
Equivalently,
\begin{equation}
  \frac{1}{\alpha}\le\frac{T_l}{T_r}\le\alpha.
\end{equation}
Inside this band the guide sticks and the adjacent rest lengths are unchanged.
Outside the band, rest length is transferred between the two spans while
preserving
\begin{equation}
  r_l' + r_r' = r_l+r_r.
  \label{eq:rest-conservation}
\end{equation}

\subsection{Left Side Too Tight}

If $T_l>\alpha T_r$, the active capstan boundary after projection is
\begin{equation}
  \frac{e_l'}{c_l} = \alpha \frac{e_r'}{c_r}.
  \label{eq:left-boundary}
\end{equation}
Conservation means
\[
  r_l'=r_l+\delta,\qquad r_r'=r_r-\delta,
\]
so the taut extensions become
\[
  e_l'=e_l-\delta,\qquad e_r'=e_r+\delta.
\]
Substituting these directly into \eqref{eq:left-boundary} gives
\[
  \frac{e_l-\delta}{c_l}
  =
  \alpha\frac{e_r+\delta}{c_r}.
\]
Solving for the rest-length transfer,
\begin{equation}
  \delta_l =
  \frac{c_r e_l-\alpha c_l e_r}{c_r+\alpha c_l}.
  \label{eq:left-delta}
\end{equation}
The implementation clamps
\[
  \delta=\clamp(\delta_l,0,r_r-r_\mathrm{min}),
\]
then applies
\[
  r_l\leftarrow r_l+\delta,\qquad
  r_r\leftarrow r_r-\delta.
\]

\subsection{Right Side Too Tight}

If $T_r>\alpha T_l$, the active boundary is
\begin{equation}
  \frac{e_r'}{c_r} = \alpha \frac{e_l'}{c_l}.
  \label{eq:right-boundary}
\end{equation}
Now
\[
  r_l'=r_l-\delta,\qquad r_r'=r_r+\delta,
\]
and therefore
\[
  e_l'=e_l+\delta,\qquad e_r'=e_r-\delta.
\]
Substitution into \eqref{eq:right-boundary} gives
\[
  \frac{e_r-\delta}{c_r}
  =
  \alpha\frac{e_l+\delta}{c_l},
\]
so
\begin{equation}
  \delta_r =
  \frac{c_l e_r-\alpha c_r e_l}{c_l+\alpha c_r}.
  \label{eq:right-delta}
\end{equation}
The implementation clamps
\[
  \delta=\clamp(\delta_r,0,r_l-r_\mathrm{min}),
\]
then applies
\[
  r_l\leftarrow r_l-\delta,\qquad
  r_r\leftarrow r_r+\delta.
\]

Equations \eqref{eq:left-delta} and \eqref{eq:right-delta} are the desired
rest-length distribution derived directly from the capstan equation.  The
capstan equation supplies the active boundary; rest-length conservation supplies
the update structure.

\section{Spin Coupling from the Same Capstan Projection}

The capstan angle $\theta$ in \eqref{eq:capstan-band} is the geometric wrap or
bend angle.  It is not the pulley spin angle.  For a circular pulley, pure spin
about the axle does not change the tangent points and therefore does not change
$\theta$.  Pulley spin enters through surface travel.  Let $\phi$ be the pulley
spin angle, and choose the sign convention that a positive scalar transfer
$\delta$ moves cable material from the left free span to the right free span:
\[
  r_l' = r_l-\delta,\qquad r_r'=r_r+\delta.
\]
The rolling no-slip candidate over one substep is then
\begin{equation}
  \delta_0 = \delta_\mathrm{geom} + R\,\dd\phi,
  \label{eq:no-slip-transfer}
\end{equation}
where $\delta_\mathrm{geom}$ is the signed tangent-point surface transfer from
the route update and $R\dd\phi$ is the rim travel caused by pulley spin.  This
candidate gives
\begin{equation}
  r_l^\star = r_l-\delta_0,\qquad
  r_r^\star = r_r+\delta_0,
  \label{eq:no-slip-rest-candidate}
\end{equation}
and therefore
\begin{equation}
  e_l^\star=\max(L_l-r_l^\star,0),\qquad
  e_r^\star=\max(L_r-r_r^\star,0),
\end{equation}
\begin{equation}
  T_l^\star=\frac{e_l^\star}{c_l},\qquad
  T_r^\star=\frac{e_r^\star}{c_r}.
\end{equation}
The no-slip rolling update is statically admissible exactly when the candidate
tensions satisfy the capstan band:
\begin{equation}
  \frac{1}{\alpha}
  \le
  \frac{T_l^\star}{T_r^\star}
  \le
  \alpha.
  \label{eq:no-slip-admissible}
\end{equation}
If \eqref{eq:no-slip-admissible} holds, static friction can support the
candidate and $\delta=\delta_0$.  If it does not hold, slip occurs and
$\delta$ is projected back to the active capstan boundary using
\eqref{eq:left-delta} or \eqref{eq:right-delta}.  In this sense, the
rest-length update and the rolling no-slip/slip decision are both obtained from
the same capstan equation: the spin angle only supplies the no-slip transfer
candidate through $R\dd\phi$.

The corresponding angular form follows by rewriting the capstan band as a
bound on the tension difference.  From $T_l=\alpha T_r$,
\[
  \frac{T_l-T_r}{T_l+T_r}
  =
  \frac{\alpha-1}{\alpha+1},
\]
so the symmetric capstan cone is
\begin{equation}
  |T_r-T_l|
  \le
  \beta(T_r+T_l),
  \qquad
  \beta=\frac{\alpha-1}{\alpha+1}.
  \label{eq:capstan-difference-cone}
\end{equation}
The tangential force imbalance on a circular pulley is $T_r-T_l$, so the
capstan-limited spin torque is
\begin{equation}
  \tau_\phi
  =
  R\,\clamp\!\left(
    T_r-T_l,\,
    -\beta(T_l+T_r),\,
    \beta(T_l+T_r)
  \right).
  \label{eq:capstan-torque}
\end{equation}
For a scalar spin inertia $I_\phi$, the spin equation of motion is
\begin{equation}
  I_\phi\dot{\omega}_\phi = \tau_\phi + \tau_\mathrm{ext},
  \label{eq:spin-eom}
\end{equation}
and a simple discrete angular impulse update is
\begin{equation}
  \dd\omega_\phi
  =
  h I_\phi^{-1}\tau_\phi.
  \label{eq:spin-velocity-update}
\end{equation}
Equations \eqref{eq:no-slip-transfer}--\eqref{eq:spin-velocity-update} are a
useful target formulation: capstan friction chooses between the no-slip
surface-transfer candidate and an active-boundary slip projection, while the
same capstan cone supplies the admissible pulley spin torque.

\section{Rolling Spin-Axis Row}

The stretch row intentionally removed pulley spin coupling in
\eqref{eq:spin-projection}.  A rolling link therefore receives a separate
angular correction, admitted only to the extent allowed by the same capstan
cone.  The current implementation uses an impulse-style approximation to
\eqref{eq:capstan-torque}: it reconstructs the spin-axis correction that the
neighboring stretch rows would have applied, then gates that correction by the
difference cone \eqref{eq:capstan-difference-cone}.

The symmetric difference/sum form of the capstan band is
\begin{equation}
  |T_l-T_r| \le \beta(T_l+T_r),\qquad
  \beta=\frac{\alpha-1}{\alpha+1}.
  \label{eq:symmetric-capstan}
\end{equation}
This is the same cone as \eqref{eq:capstan-difference-cone}.

For a rolling guide body with axis $a$, define
\begin{equation}
  s_g =
  \min\left(
    1,\;
    \frac{\beta(T_l+T_r)}{\max(|T_l-T_r|,\eps)}
  \right).
  \label{eq:spin-scale}
\end{equation}
If the adjacent stretch rows had kept the spin-axis term, their angular
corrections on the pulley would have contained
\begin{align}
  A_l &= a a^\T \left((x_l^+-x_g^c)\times n_l\right)
        \dd\lambda_l \rho,\\
  A_r &= a a^\T \left(-(x_r^--x_g^c)\times n_r\right)
        \dd\lambda_r \rho,
  \label{eq:spin-candidates}
\end{align}
where $\rho$ is the relaxation factor, $x_l^+$ is the left span endpoint on the
pulley, $x_r^-$ is the right span endpoint on the pulley, $x_g^c$ is the pulley
body center of mass, and $n_l,n_r$ are the left and right span directions.

The implemented rolling spin row applies the pure angular body delta
\begin{equation}
  \dd L_g^\mathrm{spin}
  =
  \beta\,s_g\,(A_l+A_r),
  \qquad
  \dd P_g^\mathrm{spin}=0.
  \label{eq:spin-row}
\end{equation}
This row has no independent multiplier in the current implementation.  It uses
the lagged stretch increments $\dd\lambda_l,\dd\lambda_r$ as the candidate
no-slip rolling correction and gates that correction by the capstan cone.  The
limiting cases are:
\[
  \mu=0 \Rightarrow \beta=0 \quad\text{(no pulley spin grip),}
\]
\[
  \mu\to\infty \Rightarrow \beta\to 1 \quad\text{(no-slip rolling limit).}
\]

\section{Hysteresis}

Because $r_l$ and $r_r$ are persistent state, the capstan projection reproduces
static capstan hysteresis without requiring distinct static and kinetic
coefficients.  Consider a fixed left endpoint, a rolling guide, and a driven
right endpoint.  During loading, $T_r$ increases.  When
\[
  T_r>\alpha T_l,
\]
equation \eqref{eq:right-delta} transfers rest length from left to right:
\[
  r_l\downarrow,\qquad r_r\uparrow.
\]
The left endpoint has not moved, but the left span has shorter rest length, so
the fixed side develops tension.  On that loading slip boundary,
\[
  T_r/T_l=\alpha.
\]
During unloading, the previous rest-length distribution remains latched while
the tensions stay inside the capstan band.  Once the opposite boundary is
violated,
\[
  T_l>\alpha T_r,
\]
equation \eqref{eq:left-delta} reverses the rest-length transfer and the
unloading branch follows
\[
  T_l/T_r=\alpha.
\]
The area between these two branches is the hysteresis band.

\section{XPBD and VBD Integration}

The XPBD tendon iteration is:
\begin{enumerate}
  \item update route contacts and rolling/pinhole rest-length transfer,
  \item solve each active free-span stretch row using \eqref{eq:xpbd-stretch},
  \item solve each rolling spin-axis row using \eqref{eq:spin-row}.
\end{enumerate}
The VBD rigid-body path calls the same route update, stretch row, and rolling
spin row inside the rigid iteration, using VBD's effective inverse mass and
inverse inertia arrays.  There is no separate VBD global length force and no
separate tendon Hessian model for routed cables; VBD regularizes the rigid
solve while reusing the same cable rows as XPBD.

\section{Implementation Map}

The equations above correspond to the current Newton implementation as follows:
\begin{center}
\begin{tabular}{>{\raggedright\arraybackslash}p{0.34\linewidth}
                >{\raggedright\arraybackslash}p{0.54\linewidth}}
\hline
Equation or mechanism & Implementation \\
\hline
Route tangent update and surface transfer &
\path{newton/_src/solvers/tendon_kernels.py}\newline
\texttt{update\_tendon\_attachments} \\
Stretch row \eqref{eq:xpbd-stretch} &
\path{newton/_src/solvers/xpbd/tendon_kernels.py}\newline
\texttt{solve\_tendon\_stretch} \\
Pinhole and rolling rest-length projection &
\path{newton/_src/solvers/tendon_kernels.py}\newline
\texttt{update\_tendon\_attachments} \\
Rolling spin row \eqref{eq:spin-row} &
\path{newton/_src/solvers/xpbd/tendon_kernels.py}\newline
\texttt{solve\_tendon\_slip} \\
VBD routed-tendon projection &
\path{newton/_src/solvers/vbd/solver_vbd.py}\newline
\texttt{\_project\_tendon\_constraints} \\
\hline
\end{tabular}
\end{center}

\section{Limitations}

This is a reduced massless cable model, not a resolved rope/contact model.  It
does not represent distributed cable mass, bending stiffness, contact patch
pressure, or geodesic rerouting over arbitrary geometry.  The pinhole and
rolling capstan laws use local bend or wrap angle as the reduced contact
measure.  Displacement-driven force benchmarks require compliance somewhere in
the system: cable compliance, actuator compliance, or a load-cell model.

\end{document}
